\subsection{The Most Stable Form Is Diffuse Tangential Activity}

\begin{figure*}[t]
\centering
\includegraphics[width=\textwidth]{figures/4134_figure3_spatial_timing_structure.png}
\caption{Spatial and timing structure of the T1 residual. Support fractions
and aggregate event-aligned profiles show that diffuse all-tangential activity
was the most stable repeated form. Near-pre timing, edge/core contrast, and
signed structure were retained as bounded or secondary patterns.}
\label{fig:spatial_timing}
\end{figure*}

The spatial decomposition showed that the most repeatable form of the residual
was diffuse tangential activity. Diffuse all-tangential activity passed in 13
of 14 tested observations, whereas near-pre timing passed in 8 of 14 tested
observations (Fig.~\ref{fig:spatial_timing}). Edge/core contrast appeared in a
majority subset but was weaker and less stable as a phase-specific
interpretation.

Signed event-type structure was also heterogeneous. The observations did not
support a single low-to-high, high-to-low, or mirror-symmetric signed law. The
spatial/timing result should therefore be read as evidence for a diffuse
tangential residual, not as evidence for a sharp universal trigger.

